\textbf{Phần 1. Trắc nghiệm lựa chọn}
\setcounterpageref{ex}{0}

\begin{ex} %[2D11V58]
Giá trị của biểu thức $350\text{ m} + 150\text{ m}$ là
\choice
{400 m}
{600 m}
{\True 500 m}
{450 m}
\loigiai{
$350 + 150 = 500$ m.
\textbf{Đáp án của câu hỏi là C}
}
\end{ex}

\begin{ex} %[2D11V58]
Có $3$ tờ $200\text{ đồng}$ và $2$ tờ $100\text{ đồng}$. Tổng số tiền là
\choice
{700 đồng}
{\True 800 đồng}
{600 đồng}
{900 đồng}
\loigiai{
$200 \times 3 + 100 \times 2 = 600 + 200 = 800$ đồng.
\textbf{Đáp án của câu hỏi là B}
}
\end{ex}

\begin{ex} %[2D11V58]
Một đoạn đường dài $1\text{ km}$. Người ta đã trải nhựa được $700\text{ m}$. Còn bao nhiêu mét đường chưa được trải nhựa?
\choice
{400 m}
{200 m}
{500 m}
{\True 300 m}
\loigiai{
$1\text{ km} = 1 000\text{ m}$. Còn lại: $1 000 - 700 = 300$ m.
\textbf{Đáp án của câu hỏi là D}
}
\end{ex}

\begin{ex} %[2D11V58]
Điền số thích hợp: $5\text{ m } 6\text{ dm} = \dots\text{ dm}$.
\choice
{\True 56 dm}
{506 dm}
{11 dm}
{65 dm}
\loigiai{
$5\text{ m} = 50\text{ dm} \Rightarrow 5\text{ m } 6\text{ dm} = 56$ dm.
\textbf{Đáp án của câu hỏi là A}
}
\end{ex}

\begin{ex} %[2D11V58]
Một người có $1 000\text{ đồng}$, mua một chiếc kẹo giá $300\text{ đồng}$ và một gói bim bim giá $500\text{ đồng}$. Người đó còn lại bao nhiêu tiền?
\choice
{100 đồng}
{300 đồng}
{\True 200 đồng}
{400 đồng}
\loigiai{
Đã tiêu: $300 + 500 = 800$ đồng. Còn lại: $1 000 - 800 = 200$ đồng.
\textbf{Đáp án của câu hỏi là C}
}
\end{ex}

\begin{ex} %[2D11V58]
Sợi dây thứ nhất dài $12\text{ dm}$, sợi dây thứ hai dài gấp đôi sợi thứ nhất. Cả hai sợi dây dài bao nhiêu đề-xi-mét?
\choice
{24 dm}
{\True 36 dm}
{30 dm}
{48 dm}
\loigiai{
Sợi 2: $12 \times 2 = 24$ dm. Cả hai sợi: $12 + 24 = 36$ dm.
\textbf{Đáp án của câu hỏi là B}
}
\end{ex}

\begin{ex} %[2D11V58]
Một hình chữ nhật có chiều dài $30\text{ m}$, chiều rộng $20\text{ m}$. Chu vi hình chữ nhật đó là
\choice
{50 m}
{80 m}
{60 m}
{\True 100 m}
\loigiai{
$(30 + 20) \times 2 = 100$ m.
\textbf{Đáp án của câu hỏi là D}
}
\end{ex}

\begin{ex} %[2D11V58]
Nam mua một hộp sữa giá $700\text{ đồng}$. Nam đưa cho cô bán hàng $4$ tờ tiền $200\text{ đồng}$. Cô bán hàng trả lại Nam bao nhiêu tiền?
\choice
{\True 100 đồng}
{200 đồng}
{50 đồng}
{150 đồng}
\loigiai{
Nam đưa: $200 \times 4 = 800$ đồng. Trả lại: $800 - 700 = 100$ đồng.
\textbf{Đáp án của câu hỏi là A}
}
\end{ex}

\begin{ex} %[2D11V58]
Độ dài đường gấp khúc gồm $4$ đoạn thẳng bằng nhau, mỗi đoạn dài $25\text{ cm}$ là
\choice
{10 cm}
{50 cm}
{\True 1 m}
{2 m}
\loigiai{
$25 \times 4 = 100\text{ cm} = 1\text{ m}$.
\textbf{Đáp án của câu hỏi là C}
}
\end{ex}

\begin{ex} %[2D11V58]
Điền dấu thích hợp: $1\text{ km} - 400\text{ m} \dots 500\text{ m}$.
\choice
{$<$}
{\True $>$}
{$=$}
{Không so sánh được}
\loigiai{
$1\text{ km} - 400\text{ m} = 600\text{ m} > 500\text{ m}$.
\textbf{Đáp án của câu hỏi là B}
}
\end{ex}

\begin{ex} %[2D11V58]
Bác thợ rào một khu đất hình vuông cạnh $20\text{ m}$, có để một lối đi rộng $2\text{ m}$ không rào. Chiều dài hàng rào là
\choice
{80 m}
{76 m}
{82 m}
{\True 78 m}
\loigiai{
Chu vi khu đất: $20 \times 4 = 80$ m. Hàng rào: $80 - 2 = 78$ m.
\textbf{Đáp án của câu hỏi là D}
}
\end{ex}

\begin{ex} %[2D11V58]
Lan có $1 000\text{ đồng}$, mua $5$ quả mận giá $200\text{ đồng}$ một quả. Lan còn lại bao nhiêu đồng?
\choice
{\True 0 đồng}
{200 đồng}
{100 đồng}
{500 đồng}
\loigiai{
Lan mua hết: $200 \times 5 = 1 000$ đồng. Còn lại: $1 000 - 1 000 = 0$ đồng.
\textbf{Đáp án của câu hỏi là A}
}
\end{ex}

\textbf{Phần 2. Trắc nghiệm đúng sai}
\setcounterpageref{ex}{0}

\begin{ex} %[2D11V58]
Xét tính đúng sai của các phép toán đo lường và tiền tệ:
\choiceTF
{\True 1 km dài hơn 1 000 dm.}
{3 tờ 200 đồng có giá trị ít hơn 1 tờ 500 đồng.}
{\True Một hình vuông có chu vi 1 m thì cạnh của nó dài 25 cm.}
{\True 1 000 đồng đổi được đúng 10 tờ 100 đồng.}
\loigiai{
\begin{itemize}
    \item Ý a) Đúng: 1 km = 10 000 dm > 1 000 dm. Đúng.
    \item Ý b) Sai: 200 x 3 = 600 đồng > 500 đồng.
    \item Ý c) Đúng: 100 cm : 4 = 25 cm. Đúng.
    \item Ý d) Đúng: 100 x 10 = 1 000 đồng. Đúng.
\end{itemize}
\begin{itemize}
    \item a) \textbf{Đúng}.
    \item b) \textbf{Sai}.
    \item c) \textbf{Đúng}.
    \item d) \textbf{Đúng}.
\end{itemize}
}
\end{ex}

\begin{ex} %[2D11V58]
Một cửa hàng có cuộn vải dài 100 m. Ngày thứ nhất bán được 35 m, ngày thứ hai bán được 45 m:
\choiceTF
{\True Cả hai ngày bán được 80 m vải.}
{\True Sau hai ngày, cuộn vải còn lại 20 m.}
{Số mét vải ngày thứ hai bán gấp đôi ngày thứ nhất.}
{\True Nếu mỗi mét vải giá 10 đồng thì 20 m còn lại giá 200 đồng.}
\loigiai{
\begin{itemize}
    \item Ý a) Đúng: 35 + 45 = 80 m. Đúng.
    \item Ý b) Đúng: 100 - 80 = 20 m. Đúng.
    \item Ý c) Sai: 35 x 2 = 70 != 45 m.
    \item Ý d) Đúng: 20 x 10 = 200 đồng. Đúng.
\end{itemize}
\begin{itemize}
    \item a) \textbf{Đúng}.
    \item b) \textbf{Đúng}.
    \item c) \textbf{Sai}.
    \item d) \textbf{Đúng}.
\end{itemize}
}
\end{ex}

\textbf{Phần 3. Trả lời ngắn}
\setcounterpageref{ex}{0}

\begin{ex} %[2D11V58]
Đường từ nhà An đến trường dài 400 m, từ trường đến thư viện dài 600 m. An đi từ nhà đến trường rồi qua thư viện, sau đó về nhà theo đường cũ. An đã đi được bao nhiêu ki-lô-mét?
\shortans{2}
\loigiai{
\textbf{Đáp số:} $2$\\
Một lượt: $400 + 600 = 1 000\text{ m} = 1\text{ km}$. Cả đi lẫn về: 2 km.
}
\end{ex}

\begin{ex} %[2D11V58]
Bác Hùng có 1 000 đồng, mua xà phòng hết 600 đồng, mua khăn mặt hết 300 đồng. Bác Hùng còn lại bao nhiêu đồng tiền thừa?
\shortans{100}
\loigiai{
\textbf{Đáp số:} $100$\\
Đã tiêu: $600 + 300 = 900$ đồng. Còn lại: $1 000 - 900 = 100$ đồng.
}
\end{ex}

\begin{ex} %[2D11V58]
Một sợi dây đồng dài 2 m 4 dm uốn thành hình tam giác có 3 cạnh bằng nhau. Mỗi cạnh dài bao nhiêu đề-xi-mét?
\shortans{8}
\loigiai{
\textbf{Đáp số:} $8$\\
$2\text{ m } 4\text{ dm} = 24\text{ dm}$. Mỗi cạnh: $24 : 3 = 8$ dm.
}
\end{ex}

\begin{ex} %[2D11V58]
Một cây cột cao 6 m. Người ta cắm xuống đất 1 m 5 dm. Phần cột nhô trên mặt đất dài bao nhiêu đề-xi-mét?
\shortans{45}
\loigiai{
\textbf{Đáp số:} $45$\\
6 m = 60 dm. Phần ngập đất = 15 dm. Phần nhô lên: $60 - 15 = 45$ dm.
}
\end{ex}

\begin{ex} %[2D11V58]
Mẹ mua 2 kg táo giá 500 đồng một ki-lô-gam. Mẹ đưa tờ 1 000 đồng. Người bán hàng phải trả lại mẹ bao nhiêu đồng?
\shortans{0}
\loigiai{
\textbf{Đáp số:} $0$\\
2 x 500 = 1 000 đồng. Tiền thối lại là 0 đồng.
}
\end{ex}

\begin{ex} %[2D11V58]
Một hình chữ nhật có chu vi 40 cm, chiều dài là 12 cm. Chiều rộng của hình chữ nhật đó là bao nhiêu xăng-ti-mét?
\shortans{8}
\loigiai{
\textbf{Đáp số:} $8$\\
Nửa chu vi: $40 : 2 = 20$ cm. Chiều rộng: $20 - 12 = 8$ cm.
}
\end{ex}
